\documentclass[11pt]{article}
\usepackage{mathblatt}
% Probe Serie 08.10.2026, Paket Pythagoras: Übersicht, von Hand gesetzt.
% Form wie bau/proben/2026-10-08/uebersicht-kurvenuntersuchung-3: Bereiche als Überschrift, je Eintrag ein Blatt,
% Name in Schülersprache · Fachwort, Anker nur, wo der Name nicht reicht, grau „Rand“.
% Reihenfolge: LS Kl. 9 V (1 Satz des Pythagoras, 2 in Figuren und Körpern); eigene Unterkapitel für
% Hypotenuse, Kathete, Umkehrung (Sekundo 9, Mathematik 2023 8, Fundamente 9) und Figuren/Körper.
\newcommand{\kennung}{QH4}
\newcommand{\bereich}[1]{\par\addvspace{14pt}\noindent{\large\bfseries #1}\par\nobreak\vspace{4pt}}
\newcommand{\bl}[2]{\par\noindent\makebox[1.6em][l]{$\square$}#1\ifblank{#2}{}{\hspace{0.6em}{\small\color{mbgrau}#2}}\par\vspace{3pt}}
\newcommand{\blhier}[2]{\par\noindent\llap{$\blacktriangleright$\hspace{0.5em}}\makebox[1.6em][l]{$\square$}\textbf{#1}\ifblank{#2}{}{\hspace{0.6em}{\small\color{mbgrau}#2}}\par\vspace{3pt}}
\newcommand{\rand}{\hspace{0.6em}{\small\color{mbgrau}Rand}}
\begin{document}
\pfheftstil
\fancyfoot[C]{{\footnotesize\color{mbgrau}\kennung\hfill\thepage}}
\pfheftkopf{Satz des Pythagoras}{P10 \textperiodcentered{} Übersicht}

\bereich{Vorher}
\bl{Quadrat und Wurzel \textperiodcentered{} Quadratwurzel}{$\sqrt{49} = 7$}
\bl{Rechter Winkel im Dreieck \textperiodcentered{} rechtwinkliges Dreieck}{}

\bereich{Satz des Pythagoras}
\blhier{Die lange Seite \textperiodcentered{} Hypotenuse}{$a^2 + b^2 = c^2$}
\bl{Warum der Satz stimmt \textperiodcentered{} Beweis\rand}{}
\bl{Eine kurze Seite \textperiodcentered{} Kathete}{$a^2 = c^2 - b^2$}
\bl{Ist der Winkel recht? \textperiodcentered{} Umkehrung\rand}{}

\bereich{In Figuren und Körpern}
\bl{Das Dreieck erst finden \textperiodcentered{} Pythagoras in Figuren}{Trapez, Drachen, Koordinaten}
\bl{Schräge Strecken im Körper \textperiodcentered{} Pythagoras in Körpern}{Raumdiagonale, Mantellinie}

\bereich{Weiter}
\bl{Seiten und Winkel mit sin, cos, tan \textperiodcentered{} Trigonometrie}{}
\end{document}
